% !TeX program = lualatex
% Compile twice: lualatex 224.tex
% All references and prose are embedded; no BibTeX database or external inputs.
\documentclass[11pt,a4paper]{article}
\usepackage[margin=24mm,headheight=14pt]{geometry}
\usepackage{fontspec}
\setmainfont{Latin Modern Roman}
\setsansfont{Latin Modern Sans}
\setmonofont{Latin Modern Mono}
\usepackage{amsmath,amssymb,mathtools}
\usepackage{unicode-math}
\setmathfont{Latin Modern Math}
\usepackage{microtype}
\usepackage{enumitem,xurl,needspace}
\usepackage[dvipsnames]{xcolor}
\usepackage{hyperref}
\hypersetup{unicode=true,colorlinks=true,linkcolor=MidnightBlue,urlcolor=MidnightBlue,
 pdftitle={Research Notebook: Section 224},pdfauthor={Research notebook},bookmarksnumbered=true}
\usepackage{bookmark}
\usepackage{tocloft}
\setlength{\cftsecnumwidth}{3em}
\cftsetpnumwidth{2.5em}
\cftsetrmarg{4em}
\usepackage{titlesec}
\titleformat{\section}{\Large\bfseries\raggedright}{\thesection}{0.7em}{}
\usepackage{fancyhdr}
\pagestyle{fancy}\fancyhf{}
\fancyhead[L]{\small\sffamily Open Problems | Research notebook}
\fancyhead[R]{\small 224 | 27 September 2026}
\fancyfoot[C]{\thepage}
\setlength{\parindent}{0pt}
\setlength{\parskip}{5pt plus 1pt}
\setlength{\emergencystretch}{3em}
\setcounter{tocdepth}{1}
\setcounter{secnumdepth}{1}
\setlist[itemize]{leftmargin=3em,itemsep=5pt,topsep=4pt,parsep=0pt}
\allowdisplaybreaks
\newcommand{\N}{\mathbb{N}}
\newcommand{\Z}{\mathbb{Z}}
\newcommand{\Q}{\mathbb{Q}}
\newcommand{\R}{\mathbb{R}}
\newcommand{\C}{\mathbb{C}}
\newcommand{\F}{\mathbb{F}}
\newcommand{\A}{\mathbb{A}}
\newcommand{\PP}{\mathbb P}
\newcommand{\E}{\mathbb{E}}
\newcommand{\eps}{\varepsilon}
\newcommand{\dd}{\,\mathrm d}
\newcommand{\sref}[2]{\hyperlink{source:#1:#2}{[S#2]}}
\newcommand{\eref}[2]{\hyperlink{extra:#1:#2}{[E#2]}}
% Turn the next switch off for a shorter reading copy. The quotations preserve
% every exactTarget field, including qualifications omitted from a short summary.
\newif\ifshowcatalogue
\showcataloguetrue
\newcommand{\cataloguescope}[1]{\ifshowcatalogue
 \paragraph{Exact scope in the supplied catalogue.}
 \begin{quote}\small #1\end{quote}\fi}

% Independently compilable extraction; original section text retained.
\numberwithin{equation}{section}
\begin{document}
\setcounter{section}{223}
% ================================================================
% problemId: problem.shub-smale-conjecture-for-integer-zeros
% source releaseRank: 224; source releaseStatus: open
% exactTarget SHA-256: a8c46ea7ebe430d02916344693dfe236a98748c47fa2e1c330b6988694e104fc
\clearpage
\section{\texorpdfstring{Shub-Smale \ensuremath{\tau}-Conjecture for Integer Zeros}{Shub-Smale τ-Conjecture for Integer Zeros}}\label{224}
\subsection{Definitions and mathematical statement}
A constant-free straight-line program begins with registers $1$ and $X$ and successively creates registers by addition, subtraction, or multiplication of earlier registers. Let $\tau(f)$ be the minimum number of these operations needed to output $f\in\Z[X]$. For $f\ne0$, let $Z_{\Z}(f)=\{a\in\Z:f(a)=0\}$. The conjecture asks for absolute $a>0,c\ge1$ with
\[
 |Z_{\Z}(f)|\le a(1+\tau(f))^c.
\]
Roots are counted without multiplicity, and only integer roots are counted. Arbitrary integer constants cannot be inserted free of cost; allowing that changes the complexity measure. The degree can grow exponentially in program length, so a degree bound is not the requested polynomial estimate.
\par\smallskip\textit{Formulation and concept references:} \sref{224}{1}.
\cataloguescope{Determine whether universal constants a > 0 and c >= 1 exist such that every nonzero polynomial f in Z[x] has at most a(1 + tau(f))\textasciicircum{}c distinct integer zeros, where tau(f) is the shortest constant-free straight-line-program length using addition, subtraction, and multiplication, starting from 1 and x.}
\subsection{Short English statement}
Can a polynomial produced by a short constant-free arithmetic program have more than polynomially many distinct integer roots?
% BEGIN RESEARCH ATTEMPT 224
\subsection{Research attempt}

\paragraph{Current frontier.}
Smale's Problem~4 uses the constant-free arithmetic-program measure and
records that the integer-root conjecture would imply
$\mathrm P_{\C}\ne\mathrm{NP}_{\C}$ in the Blum--Shub--Smale model.
That is a serious dependency, not a proof of Boolean $\mathrm P\ne
\mathrm{NP}$. \sref{224}{1}, Problem~4.
Rojas obtains weaker root bounds through $p$-adic additive complexity;
these are not the desired polynomial bound in unrestricted $\tau$.
\eref{224}{1}, Abstract and main theorems. Results about a \emph{real}
$\tau$ conjecture for sums of products of sparse polynomials, including
an average-case theorem, have both a different representation model and
a different root domain. \eref{224}{2}, Abstract.
The recent-source search did not supply a verified full integer-root
resolution used in this notebook.

\paragraph{Attack route.}
Degree and coefficient-height bounds are immediately available but too
large. A counterexample might instead use repeated composition to create
exponentially many preimages at small program cost, or a doubling identity
for a polynomial with many consecutive integer roots. Both mechanisms
must respect the discrete root set and the cost of substituting into an
already computed register. We test those mechanisms exactly, then isolate
what a $p$-adic branching argument would still need to prove.

\paragraph{Attempt.}
\textbf{1. Elementary bounds and what they do not buy.}
In a program of length $t$, the degree is at most $2^t$: addition does
not increase the maximum degree and multiplication at most doubles it.
Let $\|f\|_1$ be the sum of absolute values of its coefficients. Starting
from registers of norm one, each new norm is at most $2A$ or $A^2$ if
previous norms are at most $A$. Induction gives the generous bound
\[
 \|f\|_1\le 2^{2^t}.
\]
For a nonconstant nonzero integer polynomial, the usual root bound gives
$|z|\le1+\|f\|_1$ for every complex root, since the leading coefficient
has absolute value at least one. Thus only a finite, explicitly bounded
integer interval needs to be checked, but its length is doubly
exponential. These estimates supply neither a polynomial root count nor
a polynomial-time search. The zero polynomial is excluded throughout.

\textbf{2. A complete test of iterated quadratic preimages.}
For an integer $c$ set $g_c(X)=X^2-c$, and let $g_c^{\circ k}$ be its
$k$-fold iterate. Its degree is $2^k$, while its constant-free program
length is at most $\tau(c)+2k$, up to an absolute initialization constant.
We can classify all its integer zeros for every $c$ and $k\ge1$:
\begin{equation}\label{224:iterate}
 Z_{\Z}(g_c^{\circ k})=
 \begin{cases}
 \varnothing,&c<0\text{ or }c\ge0\text{ is not a square},\\
 \{0\},&c=0,\\
 \{-1,1\},&c=1,\ k\text{ odd},\\
 \{0\},&c=1,\ k\text{ even},\\
 \{-m,m\},&c=m^2,\ m\ge2,\ k=1,\\
 \varnothing,&c=m^2,\ m\ge2,\ k\ge2.
 \end{cases}
\end{equation}
Here $\tau(c)$ counts construction of the constant, rather than inserting
an arbitrary integer for free.

To prove the classification, an integer orbit ending at zero must have
its preceding integer value $y$ satisfy $y^2=c$. This excludes negative
and nonsquare $c$. For $c=m^2$, the last preimages are $\pm m$.
If $m\ge2$, their own preimages would require
$x^2=m^2-m$ or $x^2=m^2+m$. But
\[
 (m-1)^2<m^2-m<m^2<m^2+m<(m+1)^2.
\]
Neither number is a square, so no orbit of integer length at least two
can end at zero. For $c=0$ the assertion is immediate. For $c=1$,
$0\mapsto-1\mapsto0$, the preimages of zero are $\pm1$, the only
integer preimage of $-1$ is zero, and $1$ has no integer preimage.
This gives the parity rule and finishes the proof.

Thus this natural short-circuit, exponentially growing-degree family
never has more than two distinct integer roots. It is a substantive
negative test of a candidate construction, but it is only one family.
It does not control more general compositions or shared arithmetic
subcircuits. The companion script checks the complete classification
for finite ranges by exact integer preimage recursion; the displayed
argument, not that range check, proves the unrestricted family statement.

\textbf{3. A false constant-cost doubling of consecutive roots.}
Let $F_m(X)=\prod_{j=0}^{m-1}(X-j)$. There is an exact identity
\begin{equation}\label{224:falling}
 F_{2m}(X)=F_m(X)F_m(X-m).
\end{equation}
It is tempting to conclude that doubling the number of roots costs one
multiplication and one subtraction. That is not an operation allowed in
the specified model. A register containing the value $F_m(X)$ cannot be
reinterpreted as the function $F_m$ and then called at $X-m$ for free.
To use its known program at a second argument one generally must copy
and reevaluate its arithmetic gates, with the input register replaced by
$X-m$. The direct implementation therefore satisfies a bound of the form
\[
 T(2m)\le 2T(m)+O(\log m+1),
\]
where the extra cost constructs $m$, subtracts it and multiplies the
outputs. This is consistent with a linear-size program, not a logarithmic
one. It is only an upper bound for this implementation; no lower bound
excluding a different clever program for $F_m$ is proved here.

Repeated squaring does create exponentially many \emph{real} roots in
suitable Chebyshev-type recurrences. For example
$T_{2m}(X)=2T_m(X)^2-1$, with $T_1(X)=X$, yields $T_{2^k}$ in $O(k)$
constant-free operations and $2^k$ distinct real zeros in $(-1,1)$.
None is an integer: the only possible integer in that interval is zero,
and $T_{2^k}(0)$ is $\pm1$ for $k\ge1$. Thus replacing integer roots by
real roots in the desired inequality would be false, but this is not a
counterexample to the actual conjecture.

\textbf{4. A modular branching route and its remaining gap.}
Reduction modulo a prime can constrain roots, but a bound on the number
of residue roots alone does not bound integer roots: many integers can
lie in one residue class. To make the reduction injective on the actual
roots, one can take a modulus exceeding twice the root-height bound from
part~1, but that loses all polynomial control in $t$.
A more useful proposal is to bound the number of occupied branches in
the tree of residue classes modulo $p^a$, uniformly over $a$, using the
program rather than its degree. Distinct integer roots eventually split
into different branches, so such a polynomial bound would suffice.

The difficulty is at addition and subtraction gates: $u+v$ can acquire
roots where neither $u$ nor $v$ vanishes, through exact cancellation.
Multiplication only unions the two root sets, but controlling addition
by the union would be false for $u=X$ and $v=-a$, whose new root
$a\ne0$ is not a root
of either summand. More complicated cancellation can repeat at many
$p$-adic scales. The additive-complexity approach in \eref{224}{1}
quantifies some of this phenomenon, but no polynomial bound for all
constant-free programs is derived here.

\paragraph{Stress test.}
Multiplicity is never counted in \eqref{224:iterate}. The case $c=0$
therefore has one root, not $2^k$ roots. Negative constants are permitted
but must be computed; nonzero constant output polynomials have no roots.
Equation \eqref{224:falling} is correct, while the proposed gate-cost
inference is not. A failure of one efficient implementation is not a
circuit lower bound. Likewise, bounds for structured sparse circuits or
Gaussian random coefficients do not handle every integer program.

Any lemma closing the universal gap must survive the known consequence
$\mathrm P_{\C}\ne\mathrm{NP}_{\C}$ stated in the formulation source.
No step here establishes that consequence or assumes it as an input.

\paragraph{Outcome.}
\textbf{Outcome: Exploratory approach with unresolved gap.}

The attempt exactly classifies integer roots in a natural exponentially
high-degree, short-program family and diagnoses a false constant-cost
substitution argument for consecutive roots. It also separates finite
height control from the required root-count control. These family
calculations are not an improved general $\tau$ bound or a novelty claim.

The remaining problem is to control cancellation-generated integer-root
branches for arbitrary arithmetic circuits, or construct a different
short program whose \emph{distinct integer} root count grows
superpolynomially. Neither is achieved here.
\needspace{20\baselineskip}
% END RESEARCH ATTEMPT 224
\subsection{Sources}
\begin{itemize}\raggedright
\item[\textnormal{[S1]}]\hypertarget{source:224:1}{} Steve Smale, Mathematical Problems for the Next Century, 1998.
\url{https://www.fim.uni-passau.de/fileadmin/dokumente/fakultaeten/fim/lehrstuhl/muller/SmaleProblems1998.pdf}.
{\small\textit{Catalogue formulation source.}}
% BEGIN ADDED REFERENCES 224
\item[\textnormal{[E1]}]\hypertarget{extra:224:1}{}
J. M. Rojas, \emph{A Direct Ultrametric Approach to Additive Complexity
and the Shub--Smale Tau Conjecture}, arXiv:math/0304100v2, 9 April 2003,
Abstract and main root-count theorems.
\url{https://arxiv.org/abs/math/0304100v2}.
\item[\textnormal{[E2]}]\hypertarget{extra:224:2}{}
I. Briquel and P. B\"urgisser, \emph{The real tau-conjecture is true on
average}, Random Structures \& Algorithms \textbf{57} (2020), 279--303;
corrected manuscript arXiv:1806.00417v2, 6 November 2019, Abstract and
main theorem. \url{https://arxiv.org/abs/1806.00417v2}.
Both sources consulted 27 September 2026. The latter's average-case,
real-root and structured-representation hypotheses are not transferred
to the present integer-root target.
% END ADDED REFERENCES 224
\end{itemize}

\end{document}
